%% kapa: exam=1 round=37 year=2026 period=1 form=A subject=economics subject_ko=경제학원론 seq=7 number=47
%% source: Q-Net 감정평가사 기출문제 (cst00309) artlSeq=5258924 file=raw/1008/5258924/2026년 제37회 감정평가사 1차 1교시.pdf
%% conversion: kordoc 4.24.51 -> tools/gen1.py (auto); review=proofread-sonnet; flags=none
\begin{question}{47}
개인 A와 B는 하루 8시간 일하고, 쌀을 생산하거나, 밀을 생산하거나, 둘 다 생산한다. 쌀과 밀의 1kg 생산에 필요한 시간은 아래와 같다. 두 사람은 자급자족하지 않고 각자 비교우위가 있는 재화에 특화생산한 후 거래함으로써 이익을 얻고자 한다. 두 사람이 모두 거래를 통해 이익을 누리기 위하여, 쌀 1kg과 거래되는 밀의 양(kg)으로 옳은 것은?
\par
{\small\centering
\begin{tabular}{|c|c|c|}\hline
 & 쌀 1kg & 밀 1kg \\ \hline
A & 60분 & 30분 \\ \hline
B & 30분 & 10분 \\ \hline
\end{tabular}\par}
\begin{choices}
  \item 0.5
  \item 1.5
  \item 2.5
  \item 3.5
  \item 4.5
\end{choices}
\end{question}
\answer{3}
